% =====================================================================
% Related Work
% =====================================================================
\section{Related Work}
\label{sec:related}

\paragraph{Weight-space composition assumes a shared basis.}
Task arithmetic~\citep{ilharco2023taskvectors} treats a fine-tuning delta
$\theta_{\mathrm{ft}}-\theta_{\mathrm{base}}$ as a vector that can be added,
negated, or combined across tasks; TIES-Merging~\citep{yadav2023ties} and
DARE~\citep{yu2024dare} refine the combination by trimming and rescaling the
delta to reduce interference. All three operate on deltas taken \emph{from the
same base checkpoint}: the arithmetic is meaningful only because every delta
lives in one fixed parameter basis. A version bump replaces that basis. Our
naive-transfer baseline is exactly task arithmetic across versions, i.e.\ adding
$\Delta W$ to $M_{v+1}$ unchanged, and
\citet{lin2025finetuningtransfer} show empirically that this ``diff-vector
transfer'' from Llama~3.0-8B to Llama~3.1-8B retains some, but not all, of the
fine-tuning benefit. We treat their operation as the published form of the naive
baseline and ask how much of the remaining gap a data-free correction can
close.

\paragraph{Representational similarity and stitching motivate alignability.}
A long line of work shows that independently trained networks learn
representations that agree up to a simple transformation. Model
stitching~\citep{lenc2015stitching,bansal2021stitching,csiszarik2021stitching}
inserts a learned (often linear) layer between the early part of one network
and the late part of another and finds that the composite performs well;
CKA~\citep{kornblith2019cka} quantifies the residual disagreement in a
basis-invariant way; Git Re-Basin~\citep{ainsworth2023gitrebasin} aligns
permutation symmetries so that two solutions can be merged; relative
representations~\citep{moschella2023relative} remove the nuisance
transformation by re-expressing every point against a set of anchors; and the
Platonic Representation Hypothesis~\citep{huh2024platonic} argues that larger
models converge toward a shared representation of the data. These results
justify the central modeling assumption of this paper, that two versions of the
same family are related by a per-module map that is close to linear, and
anchors are the mechanism they use to find it. None of them, however, moves a
\emph{weight-level} object such as a LoRA adapter between models: stitching
adds a layer at inference time, and re-basin merges full networks rather than
re-expressing a low-rank update. Our transplant
(Eq.~\ref{eq:transplant}) folds the stitching map into the adapter so that the
ported model has the same architecture and inference cost as the original.

\paragraph{Training-free and data-free adapter transfer.}
Several recent methods port an adapter to a new base without task data; they
differ in what they require of the adapter and in what signal they align on
(Table~\ref{tab:related}).
LoRA-X~\citep{farhadzadeh2025lorax} and Trans-PEFT~\citep{gu2025transpeft}
both transfer without retraining but only for adapters trained under a special
regime: LoRA-X constrains the adapter to the dominant singular subspace of the
source base weights so that it can be re-projected, and Trans-PEFT trains
adapters to depend on attention patterns rather than on FFN-stored knowledge.
Neither applies to the pre-existing adapters an organization already has.
Cross-LoRA~\citep{xia2025crosslora} is the closest method and the baseline we
compare against. It takes an ordinary LoRA, computes rank-truncated SVDs of the
source and target base weights $W_s\!\approx\!U_s\Sigma_sV_s^\top$ and
$W_t\!\approx\!U_t\Sigma_tV_t^\top$, maps the source bases onto the target
bases by least squares, and projects the update as
$\Delta W_t=\tilde U_s(\tilde U_s^\top\Delta W_s\tilde V_s)\tilde V_s^\top$.
The signal is purely \emph{weight-space}: no input ever passes through either
model. When the two bases share dimensions, as in a version pair, the
least-squares step is exact and the method reduces to projecting $\Delta W_s$
onto the top-$k$ singular subspaces of the target weight
(\S\ref{sec:method:crosslora}). Cross-LoRA is evaluated on cross-family and
cross-size pairs (Qwen2.5-1.5B/3B, LLaMA-3.2-3B, Gemma-2-2B) on multiple-choice
benchmarks where the trained-LoRA oracle itself improves over the base by under
one point, so the reported ``approaches retrained LoRA'' leaves open how much of
a large task gain survives transfer. Two methods that do retrain are
complementary rather than competing: ReLoRA~\citep{xu2026relora} warm-starts a
new adapter from the old one with Bayesian-optimized initialization and
scheduled regularization, and PorTAL~\citep{ramp2026portal} trains a
base-agnostic latent adapter up front and fits a per-base projection from a few
hundred labeled examples. Both need task data at port time.

\begin{table}[t]
\centering
\small
\begin{tabular}{@{}lccll@{}}
\toprule
Method & Training-free & Ordinary LoRA & Alignment signal & Pairs studied \\
\midrule
Fine-tuning Transfer~\citep{lin2025finetuningtransfer} & \checkmark & \checkmark & none (diff add) & Llama 3.0$\to$3.1 \\
LoRA-X~\citep{farhadzadeh2025lorax} & \checkmark & $\times$ & weight SVD subspace & SD1.5 / SDXL \\
Trans-PEFT~\citep{gu2025transpeft} & \checkmark & $\times$ & architectural prior & 7 bases, 12 tasks \\
Cross-LoRA~\citep{xia2025crosslora} & \checkmark & \checkmark & weight SVD subspace & cross-family \\
ReLoRA~\citep{xu2026relora} & $\times$ & \checkmark & adapter as init & version updates \\
PorTAL~\citep{ramp2026portal} & $\times$ & $\times$ & learned, $\sim$500 labels & cross-size/family \\
\midrule
Ours & \checkmark & \checkmark & \textbf{activations} on unlabeled anchors & same-family versions \\
\bottomrule
\end{tabular}
\caption{Adapter transfer across base models. ``Ordinary LoRA'' means the method
accepts an adapter trained with the standard recipe, without a transfer-aware
training constraint. Only Cross-LoRA and ours are both training-free and
applicable to existing adapters; they differ in whether the alignment is read
from the weights or from what the modules compute on real inputs.}
\label{tab:related}
\end{table}

\paragraph{Where this paper sits.}
Two distinctions organize the table. The first is \emph{version versus family}:
the training-free methods above are motivated by heterogeneous pairs (different
sizes or families), and none of them studies the same-family version bump that
is the routine operational event for adapter libraries. The second is
\emph{weight versus activation}: existing training-free alignment reads the
relation between two models off their weight matrices, whereas we read it off
their activations on a small unlabeled anchor set. The two can disagree. A
version update that changes how a module responds to real inputs without
rotating the dominant singular directions of its weight matrix is invisible to a
weight-SVD method but is exactly what an activation regression captures; and an
update that is a non-orthogonal reshaping of the activation space cannot be
expressed by projection onto an orthonormal subspace at all. Our drift study
(\S\ref{sec:drift}) measures the former on real checkpoints and our ablations
(\S\ref{sec:ablations}) isolate the latter, which is why the method is
unconstrained ridge regression rather than an orthogonal map. Knowledge
distillation~\citep{hinton2015distillation} could also port behaviour to a new
base, but it needs inference on task-like inputs and gradient steps on the
target; it is a retraining method under our definition.
